% Vietnamese translation for OI Public Library
% Source-PDF: 英文资料 ENGLISH MATERIALS/Suffix Tree with Figs.pdf
\documentclass[11pt]{article}
\usepackage{oipl-vn}

\newtheorem{theorem}{Định lý}
\newtheorem{lemma}{Bổ đề}
\newenvironment{proof}{\paragraph{Chứng minh.}}{\hfill$\dashv$\par}

\newcommand{\eps}{\epsilon}
\newcommand{\STrie}{\operatorname{STrie}}
\newcommand{\STree}{\operatorname{STree}}
\newcommand{\SA}{\operatorname{SA}}

\title{Xây dựng trực tuyến cây hậu tố}
\author{Esko Ukkonen}
\date{}

\begin{document}
\maketitle

\origtitle{On-line construction of suffix trees}
\sourcepath{English Materials/Suffix Tree with Figs.pdf}

\begin{abstract}
Bài viết trình bày một thuật toán trực tuyến để xây dựng cây hậu tố của một
xâu cho trước trong thời gian tuyến tính theo độ dài xâu. Thuật toán mới có
tính chất mong muốn là xử lý xâu từng ký hiệu một từ trái sang phải, và tại
mọi thời điểm luôn có sẵn cây hậu tố của phần xâu đã quét. Phương pháp được
phát triển như một phiên bản thời gian tuyến tính của một thuật toán rất đơn
giản cho trie hậu tố, vốn có kích thước bậc hai. Dù thuật toán trie này có
trường hợp xấu nhất bậc hai, nó vẫn có thể là một phương pháp thực tế tốt khi
xâu không quá dài. Một biến thể khác của phương pháp cũng dẫn một cách tự
nhiên tới các thuật toán đã biết để xây dựng ô-tô-mát hậu tố, hay DAWG.
\end{abstract}

\noindent\textbf{Từ khóa:} thuật toán thời gian tuyến tính, cây hậu tố, trie
hậu tố, ô-tô-mát hậu tố, DAWG.

\section{Giới thiệu}

Cây hậu tố là một cấu trúc dữ liệu giống trie biểu diễn tất cả hậu tố của một
xâu. Cây hậu tố đóng vai trò trung tâm trong nhiều thuật toán trên xâu, chẳng
hạn \cite{apostolico,galil-giancarlo,amir-farach}. Tuy nhiên, các thuật toán
xây dựng cây hậu tố thời gian tuyến tính trong tài liệu thường bị xem là khá
khó nắm bắt.

Mục đích chính của bài viết này là phát triển một cách xây dựng cây hậu tố dễ
hiểu, dựa trên một ý tưởng tự nhiên giúp hoàn thiện bức tranh về cây hậu tố.
Thuật toán mới có tính chất quan trọng là trực tuyến: nó xử lý xâu từng ký
hiệu một từ trái sang phải, và luôn duy trì cây hậu tố của phần đã quét. Thuật
toán dựa trên quan sát đơn giản rằng các hậu tố của xâu
\[
  T^i=t_1\cdots t_i
\]
có thể thu được từ các hậu tố của
\[
  T^{i-1}=t_1\cdots t_{i-1}
\]
bằng cách nối ký hiệu $t_i$ vào cuối mỗi hậu tố của $T^{i-1}$, rồi thêm hậu tố
rỗng. Các hậu tố của toàn bộ xâu $T=T^n=t_1t_2\cdots t_n$ vì vậy có thể thu
được bằng cách lần lượt mở rộng hậu tố của $T^0$ thành hậu tố của $T^1$, rồi
tiếp tục cho tới khi hậu tố của $T$ được tạo từ hậu tố của $T^{n-1}$.

Cách nhìn này khác với phương pháp của Weiner \cite{weiner}, vốn đi từ phải
sang trái và thêm các hậu tố vào cây theo thứ tự tăng dần của độ dài, bắt đầu
từ hậu tố ngắn nhất; nó cũng khác với phương pháp của McCreight
\cite{mccreight}, vốn thêm các hậu tố theo thứ tự giảm dần của độ dài. Tuy
nhiên, dù trực giác ban đầu khác nhau rõ rệt, thuật toán ở đây và thuật toán
của McCreight trong dạng cuối cùng lại khá gần nhau về mặt chức năng.

Thuật toán của chúng ta được hiểu tốt nhất như một phiên bản thời gian tuyến
tính của một thuật toán khác trong \cite{ukkonen-wood} cho trie hậu tố có kích
thước bậc hai. Thuật toán rất cơ bản đó, giống thuật toán cây vị trí trong
\cite{kempf}, được trình bày ở Mục 2. Nó không chạy trong thời gian tuyến tính:
thời gian tỷ lệ với kích thước trie hậu tố, vốn có thể là bậc hai. Tuy nhiên,
một sửa đổi khá trong suốt, được mô tả ở Mục 4, cho ta phương pháp trực tuyến
thời gian tuyến tính cho cây hậu tố, đồng thời đem lại một góc nhìn tự nhiên
giúp hiểu phép xây dựng này.

Mục 5 chỉ ra thêm rằng trie hậu tố được bổ sung các liên kết hậu tố cho ta một
đặc trưng cơ bản của ô-tô-mát hậu tố, còn gọi là đồ thị từ phi chu trình có
hướng, hay DAWG. Điều này lập tức dẫn tới một thuật toán xây dựng các ô-tô-mát
đó. Phương pháp thu được về bản chất trùng với các thuật toán đã biết trong
\cite{blumer,crochemore-transducers,crochemore-constraints}. Một lần nữa,
cách nhìn mới là tự nhiên và giúp hiểu rõ các phép xây dựng ô-tô-mát hậu tố.

\section{Xây dựng trie hậu tố}

Cho $T=t_1t_2\cdots t_n$ là một xâu trên bảng chữ cái $\Sigma$. Một xâu $x$
được gọi là xâu con của $T$ nếu $T=uxv$ với một số xâu, có thể rỗng, $u$ và
$v$. Một xâu
\[
  T_i=t_i\cdots t_n,\qquad 1\le i\le n+1,
\]
là một hậu tố của $T$; đặc biệt, $T_{n+1}=\eps$ là hậu tố rỗng. Tập mọi hậu
tố của $T$ được ký hiệu là $\sigma(T)$. Trie hậu tố của $T$ là trie biểu diễn
$\sigma(T)$.

Chính xác hơn, ta ký hiệu trie hậu tố của $T$ là
\[
  \STrie(T)=(Q\cup\{\bot\},\mathit{root},F,g,f).
\]
Đây là một ô-tô-mát hữu hạn đơn định được tăng cường: đồ thị chuyển trạng thái
của nó có dạng cây và biểu diễn trie của $\sigma(T)$, còn phần tăng cường là
hàm hậu tố $f$ và trạng thái phụ $\bot$. Tập trạng thái $Q$ của $\STrie(T)$ có
thể đặt tương ứng một-một với các xâu con của $T$. Trạng thái ứng với xâu con
$x$ được ký hiệu là $\bar{x}$.

Trạng thái đầu $\mathit{root}$ ứng với xâu rỗng $\eps$, và tập trạng thái kết
thúc $F$ ứng với $\sigma(T)$. Hàm chuyển $g$ được định nghĩa bởi
\[
  g(\bar{x},a)=\bar{y}
\]
với mọi $\bar{x},\bar{y}\in Q$ sao cho $y=xa$, trong đó $a\in\Sigma$.

Hàm hậu tố $f$ được định nghĩa cho mỗi trạng thái $\bar{x}\in Q$ như sau. Nếu
$\bar{x}\ne\mathit{root}$ thì $x=ay$ với một ký hiệu $a\in\Sigma$, và ta đặt
\[
  f(\bar{x})=\bar{y}.
\]
Ngoài ra, $f(\mathit{root})=\bot$.

Trạng thái phụ $\bot$ cho phép ta viết các thuật toán phía sau mà không cần
tách riêng hậu tố rỗng và hậu tố không rỗng, hay tách riêng $\mathit{root}$
và các trạng thái khác. Trạng thái $\bot$ được nối với trie bởi
\[
  g(\bot,a)=\mathit{root}
\]
với mọi $a\in\Sigma$. Ta không định nghĩa $f(\bot)$. Có thể xem các chuyển từ
$\bot$ tới $\mathit{root}$ là nhất quán với các chuyển còn lại: trạng thái
$\bot$ ứng với nghịch đảo $a^{-1}$ của mọi ký hiệu $a\in\Sigma$; vì
$a^{-1}a=\eps$, ta đặt $g(\bot,a)=\mathit{root}$ do $\mathit{root}$ ứng với
$\eps$.

Theo thuật ngữ của McCreight \cite{mccreight}, $f(r)$ được gọi là
\term{liên kết hậu tố} của trạng thái $r$. Các liên kết hậu tố sẽ được dùng
trong quá trình xây dựng cây hậu tố; chúng cũng có nhiều ứng dụng khác
\cite{ukkonen-cpm,ukkonen-wood}.

Ô-tô-mát $\STrie(T)$ trùng với ô-tô-mát khớp xâu Aho--Corasick
\cite{aho-corasick} cho tập từ khóa $\{T_i\mid 1\le i\le n+1\}$; trong
\cite{aho-corasick}, liên kết hậu tố được gọi là chuyển thất bại.

Hình 1 của bài gốc minh họa quá trình xây dựng $\STrie(\texttt{cacao})$:
các cung chuyển trạng thái được vẽ đậm, còn các chuyển thất bại, tức liên kết
hậu tố, được vẽ mảnh. Để hình dễ đọc, chỉ hai lớp cuối của các liên kết hậu tố
được vẽ tường minh.

Dễ xây dựng $\STrie(T)$ trực tuyến bằng một lượt quét trái sang phải trên
$T$. Ký hiệu $T^i=t_1\cdots t_i$ là tiền tố của $T$, với $0\le i\le n$. Các
kết quả trung gian của phép xây dựng là $\STrie(T^i)$ với
$i=0,1,\ldots,n$.

Quan sát then chốt để nhận $\STrie(T^i)$ từ $\STrie(T^{i-1})$ là các hậu tố
của $T^i$ có thể thu được bằng cách nối $t_i$ vào cuối mỗi hậu tố của
$T^{i-1}$, rồi thêm một hậu tố rỗng:
\[
  \sigma(T^i)=\sigma(T^{i-1})t_i\cup\{\eps\}.
\]
Theo định nghĩa, $\STrie(T^{i-1})$ chấp nhận $\sigma(T^{i-1})$. Để nó chấp
nhận $\sigma(T^i)$, ta xét tập trạng thái kết thúc $F_{i-1}$ của
$\STrie(T^{i-1})$. Nếu $r\in F_{i-1}$ chưa có chuyển theo $t_i$, ta thêm một
chuyển như vậy từ $r$ tới một trạng thái mới, trạng thái này trở thành lá mới
của trie. Các trạng thái mà một chuyển cũ hoặc mới theo $t_i$ từ một trạng
thái trong $F_{i-1}$ đi tới, cùng với $\mathit{root}$, tạo thành tập trạng
thái kết thúc $F_i$ của $\STrie(T^i)$.

Các trạng thái $r\in F_{i-1}$ cần nhận chuyển mới có thể được tìm bằng liên
kết hậu tố. Định nghĩa của hàm hậu tố cho thấy $r\in F_{i-1}$ khi và chỉ khi
\[
  r=f^j(\overline{t_1\cdots t_{i-1}})
\]
với một $0\le j\le i-1$. Vì vậy mọi trạng thái trong $F_{i-1}$ nằm trên đường
đi liên kết hậu tố bắt đầu từ trạng thái sâu nhất $\overline{t_1\cdots
t_{i-1}}$ của $\STrie(T^{i-1})$ và kết thúc ở $\bot$. Ta gọi đường đi quan
trọng này là \term{đường biên} của $\STrie(T^{i-1})$.

Ta duyệt đường biên. Nếu một trạng thái $\bar{z}$ trên đường biên chưa có
chuyển theo $t_i$, ta thêm trạng thái mới $\overline{zt_i}$ và chuyển mới
\[
  g(\bar{z},t_i)=\overline{zt_i}.
\]
Như vậy $g$ được cập nhật. Để cập nhật $f$, các trạng thái mới
$\overline{zt_i}$ được nối với nhau bằng các liên kết hậu tố mới, tạo thành
một đường đi bắt đầu từ trạng thái $\overline{t_1\cdots t_i}$. Rõ ràng đây
chính là đường biên của $\STrie(T^i)$.

Việc duyệt $F_{i-1}$ dọc theo đường biên có thể dừng ngay khi gặp trạng thái
đầu tiên $\bar{z}$ sao cho trạng thái $\overline{zt_i}$, và do đó chuyển
$g(\bar{z},t_i)=\overline{zt_i}$, đã tồn tại. Thật vậy, nếu
$\overline{zt_i}$ đã là một trạng thái thì $\STrie(T^{i-1})$ phải chứa trạng
thái $\overline{z't_i}$ và chuyển $g(\bar{z'},t_i)=\overline{z't_i}$ với mọi
$z'=f^j(\bar{z})$, $j\ge 1$. Nói cách khác, nếu $zt_i$ là xâu con của
$T^{i-1}$ thì mọi hậu tố của $zt_i$ cũng là xâu con của $T^{i-1}$. Chú ý rằng
trạng thái $\bar{z}$ như vậy luôn tồn tại, vì $\bot$ là trạng thái cuối của
đường biên và $\bot$ có chuyển cho mọi ký hiệu $t_i$.

Khi phép duyệt dừng như vậy, thủ tục tạo một trạng thái mới cho mỗi liên kết
hậu tố đã xét trong lúc duyệt. Do đó toàn bộ thủ tục tốn thời gian tỷ lệ với
kích thước ô-tô-mát thu được.

\paragraph{Thuật toán 1.}
Xây dựng $\STrie(T^i)$ từ $\STrie(T^{i-1})$. Ở đây \texttt{top} là trạng thái
$\overline{t_1\cdots t_{i-1}}$.
\begin{verbatim}
r <- top;
while g(r, ti) is undefined do
    create new state r' and new transition g(r, ti) = r';
    if r != top then create new suffix link f(oldr') = r';
    oldr' <- r';
    r <- f(r);
create new suffix link f(oldr') = g(r, ti);
top <- g(top, ti).
\end{verbatim}

Bắt đầu từ $\STrie(\eps)$, gồm chỉ $\mathit{root}$, $\bot$ và các liên kết
giữa chúng, rồi lặp Thuật toán 1 cho $t_i=t_1,t_2,\ldots,t_n$, ta thu được
$\STrie(T)$. Thuật toán là tối ưu theo nghĩa thời gian chạy tỷ lệ với kích
thước kết quả cuối cùng $\STrie(T)$. Kích thước này lại tỷ lệ với $|Q|$, tức
số xâu con khác nhau của $T$. Đáng tiếc, số đó có thể là bậc hai theo $|T|$,
chẳng hạn khi $T=a^nb^n$.

\begin{theorem}
Trie hậu tố $\STrie(T)$ có thể được xây dựng trong thời gian tỷ lệ với kích
thước của $\STrie(T)$; trong trường hợp xấu nhất, kích thước này là
$O(|T|^2)$.
\end{theorem}

\section{Cây hậu tố}

Cây hậu tố $\STree(T)$ của $T$ là cấu trúc dữ liệu biểu diễn $\STrie(T)$ trong
không gian tuyến tính theo độ dài $|T|$. Điều này đạt được bằng cách chỉ biểu
diễn một tập con $Q'\cup\{\bot\}$ của các trạng thái trong $\STrie(T)$. Ta gọi
các trạng thái trong $Q'\cup\{\bot\}$ là \term{trạng thái tường minh}. Tập
$Q'$ gồm mọi trạng thái rẽ nhánh, tức các trạng thái có ít nhất hai chuyển đi
ra, và mọi lá, tức các trạng thái không có chuyển đi ra, của $\STrie(T)$. Theo
định nghĩa, $\mathit{root}$ được xem là trạng thái rẽ nhánh. Các trạng thái
còn lại của $\STrie(T)$, tức các trạng thái khác $\mathit{root}$ và $\bot$ có
đúng một chuyển đi ra, được gọi là \term{trạng thái ngầm} khi xét như trạng
thái của $\STree(T)$; chúng không xuất hiện tường minh trong $\STree(T)$.

Giả sử xâu $w$ được đọc trên đường chuyển trong $\STrie(T)$ giữa hai trạng
thái tường minh $s$ và $r$. Trong $\STree(T)$, đường đó được biểu diễn bởi một
chuyển tổng quát
\[
  g'(s,w)=r.
\]
Để tiết kiệm không gian, xâu $w$ thật ra được biểu diễn bởi một cặp con trỏ
$(k,p)$ vào $T$, gọi là con trỏ trái và con trỏ phải, sao cho
$t_k\cdots t_p=w$. Vì vậy chuyển tổng quát có dạng
\[
  g'(s,(k,p))=r.
\]

Các con trỏ như vậy tồn tại vì phải có một hậu tố $T_i$ sao cho đường chuyển
của $T_i$ trong $\STrie(T)$ đi qua $s$ và $r$. Ta có thể chọn $i$ nhỏ nhất như
vậy, rồi cho $k$ và $p$ trỏ tới đoạn con của $T_i$ được đọc trên đường từ $s$
tới $r$. Một chuyển $g'(s,(k,p))=r$ được gọi là chuyển theo $a$ nếu $t_k=a$.
Mỗi $s$ có nhiều nhất một chuyển theo $a$ với mỗi $a\in\Sigma$.

Các chuyển $g(\bot,a)=\mathit{root}$ được biểu diễn tương tự. Nếu
$\Sigma=\{a_1,a_2,\ldots,a_m\}$ thì $g(\bot,a_j)=\mathit{root}$ được biểu diễn
bởi
\[
  g'(\bot,(-j,-j))=\mathit{root},\qquad j=1,\ldots,m.
\]

Do đó cây hậu tố $\STree(T)$ có hai thành phần: bản thân cây và xâu $T$. Cấu
trúc này có kích thước tuyến tính theo $|T|$, vì $Q'$ có nhiều nhất $|T|$ lá,
mỗi hậu tố không rỗng cho nhiều nhất một lá, và do đó có nhiều nhất $|T|-1$
trạng thái rẽ nhánh khi $|T|>1$. Số chuyển giữa các trạng thái trong $Q'$ vì
thế nhiều nhất là $2|T|-2$, và mỗi chuyển chiếm không gian hằng số nhờ dùng
con trỏ thay vì lưu xâu tường minh. Ở đây ta giả sử mô hình RAM chuẩn, trong
đó một con trỏ chiếm không gian hằng số.

Ta lại tăng cường cấu trúc bằng hàm hậu tố $f'$, lần này chỉ định nghĩa cho
các trạng thái rẽ nhánh $\bar{x}\ne\mathit{root}$ bởi
\[
  f'(\bar{x})=\bar{y},
\]
trong đó $y$ là một trạng thái rẽ nhánh và $x=ay$ với một $a\in\Sigma$; ngoài
ra $f'(\mathit{root})=\bot$. Hàm $f'$ là xác định tốt: nếu $\bar{x}$ là trạng
thái rẽ nhánh thì $f'(\bar{x})$ cũng là trạng thái rẽ nhánh. Các liên kết hậu
tố này được biểu diễn tường minh. Đôi khi ta cũng sẽ nói tới liên kết hậu tố
ngầm, tức các liên kết tưởng tượng giữa các trạng thái ngầm.

Ta ký hiệu cây hậu tố của $T$ là
\[
  \STree(T)=(Q'\cup\{\bot\},\mathit{root},g',f').
\]
Một trạng thái tường minh hoặc ngầm $r$ của cây hậu tố được tham chiếu bằng
một \term{cặp tham chiếu} $(s,w)$, trong đó $s$ là một trạng thái tường minh
tổ tiên của $r$, còn $w$ là xâu đọc được trên các chuyển từ $s$ tới $r$ trong
trie hậu tố tương ứng. Cặp tham chiếu là \term{chuẩn tắc} nếu $s$ là tổ tiên
tường minh gần nhất của $r$, và vì vậy $w$ ngắn nhất có thể. Với một trạng
thái tường minh $r$, cặp tham chiếu chuẩn tắc hiển nhiên là $(r,\eps)$. Như
trên, xâu $w$ được biểu diễn bởi cặp con trỏ $(k,p)$ sao cho
$t_k\cdots t_p=w$. Nhờ vậy cặp tham chiếu $(s,w)$ có dạng $(s,(k,p))$. Cặp
$(s,\eps)$ được biểu diễn là $(s,(p+1,p))$.

Về mặt kỹ thuật, sẽ tiện hơn nếu bỏ các trạng thái kết thúc khỏi định nghĩa
cây hậu tố. Khi một ứng dụng cần trạng thái kết thúc tường minh, ta nhận chúng
miễn phí bằng cách thêm vào $T$ một ký hiệu kết thúc không xuất hiện ở nơi nào
khác trong $T$. Các lá của cây hậu tố cho xâu như vậy tương ứng một-một với
các hậu tố của $T$ và tạo thành tập trạng thái kết thúc. Một cách khác là
duyệt đường liên kết hậu tố từ lá $\bar{T}$ tới $\mathit{root}$ và làm mọi
trạng thái trên đường đó trở thành tường minh; đó chính là các trạng thái kết
thúc của $\STree(T)$. Trong nhiều ứng dụng của $\STree(T)$, vị trí bắt đầu của
mỗi hậu tố được lưu cùng trạng thái tương ứng. Cây tăng cường như vậy có thể
dùng làm chỉ mục để tìm bất kỳ xâu con nào của $T$.

\section{Xây dựng trực tuyến cây hậu tố}

Thuật toán xây dựng $\STree(T)$ sẽ mô phỏng Thuật toán 1. Phần lớn việc cần
làm là rõ ràng. Hình 2 của bài gốc minh họa các pha xây dựng
$\STree(\texttt{cacao})$; để đơn giản, hình vẽ ghi tường minh xâu gắn với mỗi
chuyển. Tuy nhiên, để nhận được thuật toán thời gian tuyến tính, một số chi
tiết cần được xem xét cẩn thận hơn.

Trước hết ta mô tả chính xác hơn Thuật toán 1 làm gì. Gọi
\[
  s_1=\overline{t_1\cdots t_{i-1}},\quad s_2,s_3,\ldots,s_i=\mathit{root},
  \quad s_{i+1}=\bot
\]
là các trạng thái của $\STrie(T^{i-1})$ trên đường biên. Gọi $j$ là chỉ số nhỏ
nhất sao cho $s_j$ không phải lá, và gọi $j'$ là chỉ số nhỏ nhất sao cho
$s_{j'}$ có chuyển theo $t_i$. Vì $s_1$ là lá còn $\bot$ là trạng thái không
phải lá có chuyển theo $t_i$, cả $j$ và $j'$ đều xác định tốt và $j\le j'$.

\begin{lemma}
Thuật toán 1 thêm vào $\STrie(T^{i-1})$ một chuyển theo $t_i$ cho mỗi trạng
thái $s_h$, $1\le h<j'$, sao cho với $1\le h<j$, chuyển mới mở rộng một nhánh
cũ của trie đang kết thúc ở lá $s_h$, còn với $j\le h<j'$, chuyển mới khởi đầu
một nhánh mới từ $s_h$. Thuật toán 1 không tạo chuyển nào khác.
\end{lemma}

Ta gọi trạng thái $s_j$ là \term{điểm hoạt động} và $s_{j'}$ là
\term{điểm kết thúc} của $\STrie(T^{i-1})$. Các trạng thái này cũng hiện diện,
tường minh hoặc ngầm, trong $\STree(T^{i-1})$. Ví dụ, trong ba cây cuối của
Hình 2 bài gốc cho xâu \texttt{cacao}, các điểm hoạt động lần lượt là
$(\mathit{root},\texttt{c})$, $(\mathit{root},\texttt{ca})$ và
$(\mathit{root},\eps)$.

Bổ đề 1 nói rằng Thuật toán 1 chèn hai nhóm chuyển theo $t_i$ vào
$\STrie(T^{i-1})$:
\begin{enumerate}
  \item Các trạng thái trên đường biên đứng trước điểm hoạt động $s_j$ nhận
  một chuyển. Chúng là lá, nên mỗi chuyển như vậy chỉ mở rộng một nhánh đã có
  của trie.
  \item Các trạng thái từ điểm hoạt động $s_j$ tới điểm kết thúc $s_{j'}$,
  không kể điểm kết thúc, nhận một chuyển mới. Chúng không phải lá, nên mỗi
  chuyển mới khởi đầu một nhánh mới.
\end{enumerate}

Ta diễn giải điều này trên cây hậu tố $\STree(T^{i-1})$. Nhóm chuyển thứ nhất,
những chuyển mở rộng nhánh đã có, có thể được cài đặt bằng cách cập nhật con
trỏ phải của mỗi chuyển biểu diễn nhánh đó. Giả sử
\[
  g'(s,(k,i-1))=r
\]
là một chuyển như vậy. Con trỏ phải phải trỏ tới vị trí cuối $i-1$ của
$T^{i-1}$, vì $r$ là lá và do đó đường đi tới $r$ phải đọc ra một hậu tố của
$T^{i-1}$ không xuất hiện ở nơi khác trong $T^{i-1}$. Chuyển sau cập nhật là
$g'(s,(k,i))=r$. Việc này chỉ làm xâu trên chuyển dài hơn, không thay đổi các
trạng thái $s$ và $r$. Nếu cập nhật tường minh mọi chuyển như vậy thì tốn quá
nhiều thời gian, nên ta dùng mẹo sau.

Mọi chuyển của $\STree(T^{i-1})$ đi tới một lá được gọi là \term{chuyển mở}.
Chuyển như vậy có dạng $g'(s,(k,i-1))=r$, trong đó con trỏ phải buộc phải trỏ
tới vị trí cuối $i-1$ của $T^{i-1}$. Vì vậy không cần biểu diễn giá trị thật
của con trỏ phải. Thay vào đó, chuyển mở được biểu diễn là
\[
  g'(s,(k,\infty))=r,
\]
trong đó $\infty$ cho biết chuyển này ``mở để lớn lên''. Thực chất,
$g'(s,(k,\infty))=r$ biểu diễn một nhánh có độ dài bất kỳ từ trạng thái $s$
tới trạng thái tưởng tượng $r$ ở vô cực. Khi chèn $t_i$ vào nhánh này, ta
không cần cập nhật tường minh con trỏ phải. Sau khi hoàn tất $\STree(T)$, các
ký hiệu $\infty$ có thể được thay bằng $n=|T|$. Như vậy nhóm chuyển thứ nhất
được cài đặt mà không gây thay đổi tường minh nào trên $\STree(T^{i-1})$.

Còn lại là cách thêm nhóm chuyển thứ hai vào $\STree(T^{i-1})$. Chúng tạo các
nhánh hoàn toàn mới bắt đầu từ các trạng thái $s_h$, $j\le h<j'$. Việc tìm
các trạng thái $s_h$ cần cẩn thận, vì ở thời điểm đó chúng không nhất thiết là
trạng thái tường minh. Chúng sẽ được tìm dọc đường biên của $\STree(T^{i-1})$
bằng cặp tham chiếu và liên kết hậu tố.

Đặt $h=j$ và giả sử $(s,w)$ là cặp tham chiếu chuẩn tắc cho $s_h$, tức điểm
hoạt động. Vì $s_h$ nằm trên đường biên của $\STrie(T^{i-1})$, $w$ phải là
một hậu tố của $T^{i-1}$. Do đó
\[
  (s,w)=(s,(k,i-1))
\]
với một $k\le i$.

Ta muốn tạo một nhánh mới bắt đầu từ trạng thái được biểu diễn bởi
$(s,(k,i-1))$. Tuy nhiên, trước hết phải kiểm tra liệu $(s,(k,i-1))$ đã tham
chiếu tới điểm kết thúc $s_{j'}$ hay chưa. Nếu có thì đã xong. Nếu chưa, cần
tạo một nhánh mới. Khi đó trạng thái $s_h$ được tham chiếu bởi $(s,(k,i-1))$
phải trở thành tường minh. Nếu nó chưa tường minh, ta tạo một trạng thái
tường minh, vẫn ký hiệu là $s_h$, bằng cách tách chuyển chứa trạng thái ngầm
tương ứng. Sau đó tạo một chuyển theo $t_i$ từ $s_h$. Chuyển này phải là
chuyển mở
\[
  g'(s_h,(i,\infty))=s'_h,
\]
trong đó $s'_h$ là một lá mới. Ngoài ra, liên kết hậu tố $f'(s_h)$ được thêm
nếu $s_h$ vừa được tạo bằng cách tách chuyển.

Tiếp theo phép xây dựng chuyển sang $s_{h+1}$. Nếu cặp tham chiếu cho $s_h$ là
$(s,(k,i-1))$, thì cặp tham chiếu chuẩn tắc cho $s_{h+1}$ là
\[
  \operatorname{canonize}(f'(s),(k,i-1)).
\]
Ở đây \textsc{canonize} làm cặp tham chiếu trở thành chuẩn tắc bằng cách cập
nhật trạng thái và con trỏ trái; con trỏ phải $i-1$ không đổi trong quá trình
chuẩn tắc hóa. Các thao tác trên được lặp lại cho $s_{h+1}$, rồi tiếp tục cho
tới khi tìm thấy điểm kết thúc $s_{j'}$.

Ta thu được thủ tục \textsc{update} dưới đây, biến $\STree(T^{i-1})$ thành
$\STree(T^i)$ bằng cách chèn các chuyển theo $t_i$ thuộc nhóm thứ hai. Thủ tục
dùng \textsc{canonize} vừa nêu, và thủ tục \textsc{test-and-split} để kiểm tra
một cặp tham chiếu có chỉ tới điểm kết thúc hay không. Nếu không, thủ tục tạo
và trả về một trạng thái tường minh cho cặp tham chiếu đó, trừ khi cặp đã biểu
diễn một trạng thái tường minh. Thủ tục \textsc{update} trả về một cặp tham
chiếu cho điểm kết thúc $s_{j'}$; trên thực tế chỉ trả về trạng thái và con
trỏ trái, vì con trỏ phải vẫn là $i-1$ cho mọi trạng thái trên đường biên.

\paragraph{Thủ tục \textsc{update}$(s,(k,i))$.}
Giả sử $(s,(k,i-1))$ là cặp tham chiếu chuẩn tắc cho điểm hoạt động.
\begin{verbatim}
oldr <- root;
(end-point, r) <- test-and-split(s, (k, i - 1), ti);
while not end-point do
    create new transition g'(r, (i, infinity)) = r',
        where r' is a new state;
    if oldr != root then create new suffix link f'(oldr) = r;
    oldr <- r;
    (s, k) <- canonize(f'(s), (k, i - 1));
    (end-point, r) <- test-and-split(s, (k, i - 1), ti);
if oldr != root then create new suffix link f'(oldr) = s;
return (s, k).
\end{verbatim}

Thủ tục \textsc{test-and-split} kiểm tra trạng thái có cặp tham chiếu chuẩn
tắc $(s,(k,p))$ có phải điểm kết thúc hay không, tức trong $\STrie(T^{i-1})$
nó đã có chuyển theo $t_i$ hay chưa. Ký hiệu $t_i$ được truyền vào dưới tên
$t$. Kết quả kiểm tra được trả về ở tham số thứ nhất. Nếu $(s,(k,p))$ không
phải điểm kết thúc, trạng thái đó được làm tường minh, nếu cần bằng cách tách
một chuyển, và trạng thái tường minh được trả về ở tham số thứ hai.

\paragraph{Thủ tục \textsc{test-and-split}$(s,(k,p),t)$.}
\begin{verbatim}
if k <= p then
    let g'(s, (k', p')) = s' be the tk-transition from s;
    if t = t_{k' + p - k + 1} then return (true, s)
    else
        replace the tk-transition above by transitions
            g'(s, (k', k' + p - k)) = r
            g'(r, (k' + p - k + 1, p')) = s',
        where r is a new state;
        return (false, r)
else
    if there is no t-transition from s then return (false, s)
    else return (true, s).
\end{verbatim}

Thủ tục này tận dụng việc $(s,(k,p))$ là chuẩn tắc: câu trả lời cho phép kiểm
tra điểm kết thúc có thể tìm trong thời gian hằng số bằng cách chỉ xét một
chuyển từ $s$.

Thủ tục \textsc{canonize} như sau. Với một cặp tham chiếu $(s,(k,p))$ cho một
trạng thái $r$, thủ tục tìm và trả về trạng thái $s'$ cùng con trỏ trái $k'$
sao cho $(s',(k',p))$ là cặp tham chiếu chuẩn tắc cho $r$. Trạng thái $s'$ là
tổ tiên tường minh gần nhất của $r$, hoặc chính $r$ nếu $r$ tường minh. Vì
vậy xâu đi từ $s'$ tới $r$ phải là một hậu tố của xâu $t_k\cdots t_p$ đi từ
$s$ tới $r$. Do đó con trỏ phải $p$ không đổi, còn con trỏ trái $k$ có thể
trở thành $k'\ge k$.

\paragraph{Thủ tục \textsc{canonize}$(s,(k,p))$.}
\begin{verbatim}
if p < k then return (s, k)
else
    find the tk-transition g'(s, (k', p')) = s' from s;
    while p' - k' <= p - k do
        k <- k + p' - k' + 1;
        s <- s';
        if k <= p then
            find the tk-transition g'(s, (k', p')) = s' from s;
    return (s, k).
\end{verbatim}

Để tiếp tục xây dựng cho ký hiệu văn bản kế tiếp $t_{i+1}$, cần tìm điểm hoạt
động của $\STree(T^i)$. Trước hết, $s_j$ là điểm hoạt động của
$\STree(T^{i-1})$ khi và chỉ khi $s_j=\overline{t_j\cdots t_{i-1}}$, trong đó
$t_j\cdots t_{i-1}$ là hậu tố dài nhất của $T^{i-1}$ xuất hiện ít nhất hai
lần trong $T^{i-1}$. Tiếp theo, $s_{j'}$ là điểm kết thúc của
$\STree(T^{i-1})$ khi và chỉ khi $s_{j'}=\overline{t_{j'}\cdots t_{i-1}}$,
trong đó $t_{j'}\cdots t_{i-1}$ là hậu tố dài nhất của $T^{i-1}$ sao cho
$t_{j'}\cdots t_{i-1}t_i$ là xâu con của $T^{i-1}$. Nhưng điều này nghĩa là
nếu $s_{j'}$ là điểm kết thúc của $\STree(T^{i-1})$ thì
$t_{j'}\cdots t_{i-1}t_i$ là hậu tố dài nhất của $T^i$ xuất hiện ít nhất hai
lần trong $T^i$, tức trạng thái $g(s_{j'},t_i)$ là điểm hoạt động của
$\STree(T^i)$.

\begin{lemma}
Giả sử $(s,(k,i-1))$ là một cặp tham chiếu của điểm kết thúc $s_{j'}$ của
$\STree(T^{i-1})$. Khi đó $(s,(k,i))$ là một cặp tham chiếu của điểm hoạt động
của $\STree(T^i)$.
\end{lemma}

Thuật toán tổng thể để xây dựng $\STree(T)$ như sau. Xâu $T$ được xử lý từng
ký hiệu một trong một lượt quét trái sang phải. Viết
$\Sigma=\{t_{-1},\ldots,t_{-m}\}$ giúp trình bày các chuyển từ $\bot$ giống
các chuyển khác.

\paragraph{Thuật toán 2.}
Xây dựng $\STree(T)$ cho xâu $T=t_1t_2\cdots t_n]$ trên bảng chữ cái
$\Sigma=\{t_{-1},\ldots,t_{-m}\}$, trong đó $]$ là ký hiệu kết thúc không xuất
hiện ở nơi khác trong $T$.
\begin{verbatim}
create states root and bottom;
for j <- 1, ..., m do
    create transition g'(bottom, (-j, -j)) = root;
create suffix link f'(root) = bottom;
s <- root; k <- 1; i <- 0;
while t_{i+1} != ] do
    i <- i + 1;
    (s, k) <- update(s, (k, i));
    (s, k) <- canonize(s, (k, i)).
\end{verbatim}

Bước cuối trong mỗi vòng dựa trên Bổ đề 2: sau lời gọi \textsc{update}, cặp
$(s,(k,i-1))$ tham chiếu tới điểm kết thúc của $\STree(T^{i-1})$, và do đó
$(s,(k,i))$ tham chiếu tới điểm hoạt động của $\STree(T^i)$.

\begin{theorem}
Thuật toán 2 xây dựng trực tuyến cây hậu tố $\STree(T)$ cho xâu
$T=t_1\cdots t_n$ trong thời gian $O(n)$.
\end{theorem}

\begin{proof}
Thuật toán xây dựng $\STree(T)$ qua các cây trung gian
\[
  \STree(T^0),\STree(T^1),\ldots,\STree(T^n)=\STree(T).
\]
Nó là trực tuyến vì để xây dựng $\STree(T^i)$, thuật toán chỉ cần truy cập
$i$ ký hiệu đầu tiên của $T$.

Để phân tích thời gian chạy, ta chia chi phí thành hai phần, và cả hai đều là
$O(n)$. Phần thứ nhất là tổng thời gian của thủ tục \textsc{canonize}. Phần
thứ hai gồm phần còn lại: thời gian lặp lại việc đi theo đường liên kết hậu tố
từ điểm hoạt động hiện tại tới điểm kết thúc, tạo các nhánh mới bằng
\textsc{update}, rồi tìm điểm hoạt động kế tiếp bằng cách đi một chuyển từ
điểm kết thúc, tức bước cuối của Thuật toán 2. Ta gọi các trạng thái, hay các
cặp tham chiếu, trên các đường này là \term{trạng thái được thăm}.

Phần thứ hai tốn thời gian tỷ lệ với tổng số trạng thái được thăm, vì các thao
tác tại mỗi trạng thái như tạo một trạng thái tường minh và nhánh mới, đi theo
một liên kết hậu tố tường minh hoặc ngầm, và kiểm tra điểm kết thúc đều có thể
cài đặt trong thời gian hằng số khi loại trừ \textsc{canonize}. Nói chính xác
hơn, điều này cũng cần giả thiết $|\Sigma|$ bị chặn độc lập với $n$. Gọi
$r_i$ là điểm hoạt động của $\STree(T^i)$ với $0\le i\le n$. Các trạng thái
được thăm giữa $r_{i-1}$ và $r_i$ nằm trên một đường gồm một số liên kết hậu
tố và một chuyển theo $t_i$. Đi theo một liên kết hậu tố làm giảm độ sâu, tức
độ dài xâu đọc được trên đường từ $\mathit{root}$, đi một đơn vị; còn đi theo
chuyển $t_i$ làm tăng nó một đơn vị. Do đó số trạng thái được thăm trên đường,
tính cả $r_{i-1}$ nhưng không tính $r_i$, là
\[
  \operatorname{depth}(r_{i-1})-\operatorname{depth}(r_i)+2.
\]
Tổng số trạng thái được thăm là
\[
  \sum_{i=1}^n
  \bigl(\operatorname{depth}(r_{i-1})-\operatorname{depth}(r_i)+2\bigr)
  =\operatorname{depth}(r_0)-\operatorname{depth}(r_n)+2n\le 2n.
\]
Vì vậy phần thời gian thứ hai là $O(n)$.

Thời gian của mỗi lần gọi \textsc{canonize} bị chặn bởi $a+bq$, với $a,b$ là
hằng số và $q$ là số lần thực hiện thân vòng lặp trong \textsc{canonize}. Tổng
thời gian của \textsc{canonize} vì vậy bị chặn tỷ lệ với tổng số lời gọi
\textsc{canonize} cộng với tổng số lần thực hiện thân vòng lặp trong mọi lời
gọi. Có $O(n)$ lời gọi, vì mỗi trạng thái được thăm gây ra một lời gọi, hoặc
ở trong \textsc{update}, hoặc trực tiếp ở bước cuối của Thuật toán 2. Mỗi lần
thực hiện thân vòng lặp xóa một xâu không rỗng khỏi đầu trái của xâu
$w=t_k\cdots t_p$ được biểu diễn bởi các con trỏ trong cặp tham chiếu
$(s,(k,p))$. Trong toàn bộ quá trình, $w$ chỉ có thể tăng ở bước cuối của
Thuật toán 2, nơi nối $t_i$ với $i=1,\ldots,n$ vào đầu phải của $w$. Do đó
việc xóa một đoạn không rỗng chỉ có thể xảy ra nhiều nhất $n$ lần. Tổng thời
gian cho thân vòng lặp là $O(n)$, và toàn bộ phần \textsc{canonize}, tức phần
thời gian thứ nhất, cũng là $O(n)$.
\end{proof}

\paragraph{Nhận xét 1.}
Theo J. Karkkainen, trong dạng cuối cùng, thuật toán này là họ hàng khá gần
của phương pháp McCreight \cite{mccreight}. Khác biệt kỹ thuật chính dường
như là: mỗi lần thực hiện thân vòng lặp chính của Thuật toán 2 tiêu thụ một
ký hiệu văn bản $t_i$, trong khi mỗi lần thực hiện thân vòng lặp chính của
thuật toán McCreight đi qua một liên kết hậu tố và tiêu thụ không hoặc nhiều
ký hiệu văn bản.

\paragraph{Nhận xét 2.}
Không khó tổng quát hóa Thuật toán 2 cho phiên bản động sau của bài toán cây
hậu tố, tương tự bài toán khớp từ điển thích nghi trong \cite{amir-farach}:
duy trì một cây hậu tố tổng quát có kích thước tuyến tính biểu diễn mọi hậu tố
của các xâu $T_i$ trong tập $\{T_1,\ldots,T_k\}$ dưới các phép chèn hoặc xóa
một xâu $T_i$. Thuật toán thu được thực hiện mỗi cập nhật trong thời gian
$O(|T_i|)$.

\section{Xây dựng ô-tô-mát hậu tố}

Ô-tô-mát hậu tố $\SA(T)$ của một xâu $T=t_1\cdots t_n$ là DFA tối tiểu chấp
nhận tất cả hậu tố của $T$.

Vì $\STrie(T)$ là một DFA cho các hậu tố của $T$, ta có thể thu được
$\SA(T)$ bằng cách tối tiểu hóa $\STrie(T)$ theo cách chuẩn. Tối tiểu hóa hoạt
động bằng cách gộp các trạng thái tương đương, tức các trạng thái mà từ đó
$\STrie(T)$ chấp nhận cùng một tập xâu. Dùng liên kết hậu tố, ta nhận được một
đặc trưng tự nhiên cho các trạng thái tương đương như sau.

Một trạng thái $s$ của $\STrie(T)$ được gọi là \term{thiết yếu} nếu có ít nhất
hai liên kết hậu tố khác nhau trỏ tới $s$, hoặc nếu
$s=\overline{t_1\cdots t_k}$ với một $k$ nào đó.

\begin{theorem}
Cho $s$ và $r$ là hai trạng thái của $\STrie(T)$. Tập xâu được chấp nhận từ
$s$ bằng tập xâu được chấp nhận từ $r$ khi và chỉ khi đường liên kết hậu tố
bắt đầu từ $s$ chứa $r$, hoặc đường bắt đầu từ $r$ chứa $s$, và đường con từ
$s$ tới $r$, hoặc từ $r$ tới $s$, không chứa trạng thái thiết yếu nào khác
ngoài, có thể, chính $s$ hoặc chính $r$.
\end{theorem}

\begin{proof}
Định lý suy ra từ các quan sát sau.

Tập xâu được chấp nhận từ một trạng thái của $\STrie(T)$ là một tập con của
các hậu tố của $T$, nên mỗi xâu được chấp nhận có độ dài khác nhau.

Một xâu độ dài $i$ được chấp nhận từ trạng thái $s$ của $\STrie(T)$ khi và
chỉ khi đường liên kết hậu tố bắt đầu từ trạng thái
$\overline{t_1\cdots t_{n-i}}$ chứa $s$.

Các liên kết hậu tố tạo thành một cây có hướng về gốc $\mathit{root}$.
\end{proof}

Điều này gợi ý một phương pháp xây dựng $\SA(T)$ bằng một Thuật toán 1 đã sửa
đổi. Điểm mới là phép xây dựng chỉ nên tạo một trạng thái mới nếu trạng thái
đó là thiết yếu. Một trạng thái không thiết yếu $s$ được gộp với trạng thái
thiết yếu đầu tiên đứng trước $s$ trên đường liên kết hậu tố đi qua $s$. Theo
Định lý 3, các trạng thái đó là tương đương, nên việc gộp là đúng.

Vì chỉ có $O(|T|)$ trạng thái thiết yếu, thuật toán thu được có thể được cài
đặt để chạy trong thời gian tuyến tính. Thuật toán này hóa ra tương tự các
thuật toán trong \cite{blumer,crochemore-transducers,crochemore-constraints},
nên bài viết gốc không trình bày chi tiết.

\section*{Lời cảm ơn}

J. Karkkainen đã chỉ ra một số điểm thiếu chính xác trong phiên bản trước của
công trình này \cite{ukkonen-ip92}. Tác giả cũng cảm ơn E. Sutinen, D. Wood,
và đặc biệt là S. Kurtz cùng G. A. Stephen vì nhiều nhận xét hữu ích.

\begin{thebibliography}{13}

\bibitem{aho-corasick}
A. Aho and M. Corasick.
Efficient string matching: An aid to bibliographic search.
\textit{Communications of the ACM}, 18:333--340, 1975.

\bibitem{amir-farach}
A. Amir and M. Farach.
Adaptive dictionary matching.
In \textit{Proc. 32nd IEEE Annual Symposium on Foundations of Computer
Science}, pages 760--766, 1991.

\bibitem{apostolico}
A. Apostolico.
The myriad virtues of subword trees.
In A. Apostolico and Z. Galil, editors, \textit{Combinatorial Algorithms on
Words}, pages 85--95. Springer-Verlag, 1985.

\bibitem{blumer}
A. Blumer et al.
The smallest automaton recognizing the subwords of a text.
\textit{Theoretical Computer Science}, 40:31--55, 1985.

\bibitem{crochemore-transducers}
M. Crochemore.
Transducers and repetitions.
\textit{Theoretical Computer Science}, 45:63--86, 1986.

\bibitem{crochemore-constraints}
M. Crochemore.
String matching with constraints.
In M. P. Chytil, L. Janiga, and V. Koubek, editors,
\textit{Mathematical Foundations of Computer Science 1988}, Lecture Notes in
Computer Science 324, pages 44--58. Springer-Verlag, 1988.

\bibitem{galil-giancarlo}
Z. Galil and R. Giancarlo.
Data structures and algorithms for approximate string matching.
\textit{Journal of Complexity}, 4:33--72, 1988.

\bibitem{kempf}
M. Kempf, R. Bayer, and U. Guntzer.
Time optimal left to right construction of position trees.
\textit{Acta Informatica}, 24:461--474, 1987.

\bibitem{mccreight}
E. McCreight.
A space-economical suffix tree construction algorithm.
\textit{Journal of the ACM}, 23:262--272, 1976.

\bibitem{ukkonen-ip92}
E. Ukkonen.
Constructing suffix trees on-line in linear time.
In J. van Leeuwen, editor, \textit{Algorithms, Software, Architecture:
Information Processing 92}, volume I, pages 484--492. Elsevier, 1992.

\bibitem{ukkonen-cpm}
E. Ukkonen.
Approximate string-matching over suffix trees.
In A. Apostolico, M. Crochemore, Z. Galil, and U. Manber, editors,
\textit{Combinatorial Pattern Matching, CPM'93}, Lecture Notes in Computer
Science 684, pages 228--242. Springer-Verlag, 1993.

\bibitem{ukkonen-wood}
E. Ukkonen and D. Wood.
Approximate string matching with suffix automata.
\textit{Algorithmica}, 10:353--364, 1993.

\bibitem{weiner}
P. Weiner.
Linear pattern matching algorithms.
In \textit{IEEE 14th Annual Symposium on Switching and Automata Theory},
pages 1--11, 1973.

\end{thebibliography}

\end{document}
